\section{Beyond Attention: Recurrent Mixers}

\begin{frame}{}
    \LARGE Next-Generation Architectures: \textbf{Beyond Attention}

    \vspace{1em}
    \normalsize Does the sequence mixer have to be attention at all?
\end{frame}

\subsection{Linear Attention}

\begin{frame}{Attention's Two Bills}
    \begin{itemize}\itemsep0.5em
        \item Softmax attention is paid for twice: once in \textbf{training compute}, once in \textbf{inference memory}.
    \end{itemize}
    \vspace{0.3em}
    \begin{center}
    \begin{adjustbox}{max width=\linewidth}\begin{tabular}{lcc}
        \toprule
        & \textbf{Training} & \textbf{State at inference} \\
        \midrule
        Classic RNN (LSTM) & sequential & fixed, $O(1)$ \\
        Softmax attention & parallel, $O(L^2)$ & KV cache, $O(L)$ \\
        \rowcolor{KAUSTcyan!12} What we want & parallel, $O(L)$ & fixed, $O(1)$ \\
        \bottomrule
    \end{tabular}\end{adjustbox}
    \end{center}
    \vspace{0.3em}
    \begin{itemize}\itemsep0.5em
        \item Albert Gu's picture: a Transformer is a \emph{database} that stores every token. A recurrent model is a \emph{brain} that keeps a compressed state.
        \item The database gives exact recall. The brain gives constant cost. \textbf{This section is about getting both.}
    \end{itemize}
    \src{Table after Nayebi and Gormley, CMU 10-423 Lecture 21 (2026). Gu, ``On the Tradeoffs of State Space Models and Transformers'' (2025), \url{https://goombalab.github.io/blog/2025/tradeoffs/}.}
\end{frame}

\begin{frame}{Linear Attention: Move One Bracket}
    \begin{itemize}\itemsep0.4em
        \item Replace the softmax similarity with a kernel, $\mathrm{sim}(q,k)=\phi(q)^\top\phi(k)$. The sum over keys then factors out of the query:
    \end{itemize}
    \[
        o_i \;=\; \frac{\sum_{j\leq i}\phi(q_i)^\top\phi(k_j)\,v_j}{\sum_{j\leq i}\phi(q_i)^\top\phi(k_j)}
        \;=\; \frac{\phi(q_i)^\top \overbrace{\textstyle\sum_{j\leq i}\phi(k_j)\,v_j^\top}^{\text{does not depend on } i\text{'s query}}}{\phi(q_i)^\top\sum_{j\leq i}\phi(k_j)}
    \]
    \begin{itemize}\itemsep0.4em
        \item In matrix form this is only associativity: $(QK^\top)V = Q\,(K^\top V)$.
        \item Cost falls from $O(L^2 d)$ to $O(L d^2)$. Stanford's CS336 calls the trick \emph{``very silly, but surprisingly important''}.
        \item Modern variants drop the normaliser and the feature map, and normalise the output instead.
    \end{itemize}
    \src{Katharopoulos et al., ``Transformers are RNNs,'' ICML 2020, \href{https://arxiv.org/abs/2006.16236}{arXiv:2006.16236}. Hashimoto, Stanford CS336 Lecture 4 (2026). Yang et al., \href{https://arxiv.org/abs/2406.06484}{arXiv:2406.06484}.}
\end{frame}

\begin{frame}{The Same Layer Is an RNN with a Matrix State}
    \begin{columns}[T,onlytextwidth]
        \column{0.52\textwidth}
        Read the sum causally and it becomes a recurrence:
        \[
            S_t = S_{t-1} + v_t k_t^\top, \qquad o_t = S_t\, q_t
        \]
        \begin{itemize}\itemsep0.45em
            \item $S_t \in \mathbb{R}^{d\times d}$ is a \textbf{fast-weight matrix}. Each token writes one rank-1 update into it.
            \item \textbf{Train} in the parallel form. \textbf{Decode} in the recurrent form, with no cache that grows.
            \item Every model in this section changes only two things: what multiplies $S_{t-1}$, and what gets written.
        \end{itemize}
        \column{0.45\textwidth}
        \centering
        \begin{adjustbox}{max width=\linewidth}\begin{tikzpicture}[font=\scriptsize, >=Stealth]
            % KV cache
            \node[font=\footnotesize\bfseries] at (1.5,3.0) {Softmax attention};
            \foreach \i in {0,...,5} {
                \draw[fill=KAUSTcyan!25] (0.5*\i,2.1) rectangle (0.5*\i+0.42,2.6);
            }
            \node at (3.35,2.35) {$\cdots$};
            \node[gray] at (1.5,1.8) {one $(k,v)$ per past token};
            % matrix state
            \node[font=\footnotesize\bfseries] at (1.5,1.1) {Linear attention};
            \draw[fill=darkorange!25] (0.9,-0.5) rectangle (2.1,0.7);
            \node at (1.5,0.1) {$S_t$};
            \draw[->, thick] (-0.2,0.1) -- (0.85,0.1) node[midway, above] {$v_t k_t^\top$};
            \draw[->, thick] (2.15,0.1) -- (3.2,0.1) node[midway, above] {$S_t q_t$};
            \node[gray] at (1.5,-0.8) {one $d\times d$ matrix, any length};
        \end{tikzpicture}\end{adjustbox}
    \end{columns}
    \src{Katharopoulos et al., \href{https://arxiv.org/abs/2006.16236}{arXiv:2006.16236}. Notation follows Yang et al., ``Parallelizing Linear Transformers with the Delta Rule,'' \href{https://arxiv.org/abs/2406.06484}{arXiv:2406.06484}.}
\end{frame}

\begin{frame}{Why It Lost for Years: Recall}
    \begin{itemize}\itemsep0.45em
        \item Ask the state for the value stored under key $k_j$:
    \end{itemize}
    \[
        S\,k_j \;=\; v_j \;+\; \underbrace{\textstyle\sum_{i\neq j}\,(k_i^\top k_j)\,v_i}_{\text{interference}}
    \]
    \begin{itemize}\itemsep0.45em
        \item Retrieval is exact only for orthogonal keys, and $\mathbb{R}^d$ holds only $d$ of them. Errors pile up with length.
        \item Writes are \textbf{additive}. The memory never forgets and never corrects a stale entry.
        \item The loss is concentrated: \alert{6.4\% of tokens} are recall lookups, and they explain \alert{82\%} of the perplexity gap to attention.
    \end{itemize}
    \vspace{0.2em}
    \takeaway{Two abilities are missing: \textbf{forgetting} and \textbf{overwriting}. The next two slides add them one at a time.}
    \src{Yang, ``DeltaNet Explained (Part I)'' (2024), \url{https://sustcsonglin.github.io/blog/2024/deltanet-1/}. Arora et al., ``Zoology,'' \href{https://arxiv.org/abs/2312.04927}{arXiv:2312.04927}.}
\end{frame}

\begin{frame}{Fix 1: Learn to Forget}
    \begin{itemize}\itemsep0.4em
        \item Multiply the old state by a decay before writing:
    \end{itemize}
    \[
        S_t \;=\; \textcolor{KAred}{G_t}\; S_{t-1} \;+\; v_t k_t^\top
    \]
    \begin{center}
    \small
    \begin{adjustbox}{max width=\linewidth}\begin{tabular}{lll}
        \toprule
        \textbf{Model} & \textbf{Decay $G_t$} & \textbf{Depends on the input?} \\
        \midrule
        RetNet (2023) & one fixed scalar $\gamma$ per head & no \\
        Mamba-2 (2024) & scalar $\gamma_t$ & yes \\
        Gated Linear Attention (2023) & per-channel $\mathrm{Diag}(\alpha_t)$ & yes \\
        \bottomrule
    \end{tabular}\end{adjustbox}
    \end{center}
    \begin{itemize}\itemsep0.4em
        \item Data-dependent gates close most of the language-modelling gap. At 1.3B parameters GLA reaches perplexity 17.22 against 16.85 for a tuned Transformer.
        \item Fixed decay did not survive. In a 2025 study, RetNet gives near-zero recall inside hybrids.
        \item \textbf{What a gate cannot do:} it scales the whole memory. It cannot fix one wrong association.
    \end{itemize}
    \src{Yang et al., ``Gated Linear Attention Transformers,'' \href{https://arxiv.org/abs/2312.06635}{arXiv:2312.06635}, Table 2. ``A Systematic Analysis of Hybrid Linear Attention,'' \href{https://arxiv.org/abs/2507.06457}{arXiv:2507.06457}.}
\end{frame}

\begin{frame}{Fix 2: Learn to Overwrite (the Delta Rule)}
    \begin{itemize}\itemsep0.4em
        \item Treat the state as a regression model that should map $k_t \mapsto v_t$, and take \textbf{one gradient step per token} on
        $\mathcal{L}_t(S)=\tfrac12\lVert S k_t - v_t\rVert^2$:
    \end{itemize}
    \[
        S_t \;=\; S_{t-1} - \beta_t\,(S_{t-1}k_t - v_t)\,k_t^\top
        \;=\; S_{t-1}\,\textcolor{KAred}{(I-\beta_t k_t k_t^\top)} \;+\; \beta_t\, v_t k_t^\top
    \]
    \begin{itemize}\itemsep0.4em
        \item The red factor \textbf{erases} whatever was stored along the current key. Then the new value is written.
        \item Plain linear attention is the same step on a loss that never looks at what is already stored.
        \item \textbf{DeltaNet} solves the hardest synthetic recall setting (MQAR) at 100\%.
        \item Its weakness mirrors Fix 1. With no decay, clearing the whole memory takes $d$ steps.
    \end{itemize}
    \src{Schlag et al., ``Linear Transformers Are Secretly Fast Weight Programmers,'' ICML 2021. Yang et al., ``Parallelizing Linear Transformers with the Delta Rule over Sequence Length,'' NeurIPS 2024, \href{https://arxiv.org/abs/2406.06484}{arXiv:2406.06484}.}
\end{frame}

\begin{frame}{Gated DeltaNet: Forget \emph{and} Overwrite}
    \[
        S_t \;=\; S_{t-1}\,\big(\textcolor{KAUSTblue}{\alpha_t}\,\textcolor{KAred}{(I-\beta_t k_t k_t^\top)}\big) \;+\; \beta_t\, v_t k_t^\top
    \]
    \begin{columns}[T,onlytextwidth]
        \column{0.47\textwidth}
        \begin{itemize}\itemsep0.45em
            \item \textcolor{KAUSTblue}{Gate}: fast, global erase.
            \item \textcolor{KAred}{Delta rule}: targeted edit of one key.
            \item The two are complementary, not alternatives.
            \item This is the mixer inside \textbf{Qwen3-Next} and \textbf{Qwen3.5}. Kimi's \textbf{KDA} uses a per-channel gate in place of the scalar.
        \end{itemize}
        \column{0.5\textwidth}
        \centering
        \footnotesize
        \begin{adjustbox}{max width=\linewidth}\begin{tabular}{lccc}
            \toprule
            \textbf{1.3B, 100B tokens} & \textbf{PPL} $\downarrow$ & \textbf{Recall} $\uparrow$ & \textbf{Long} $\uparrow$ \\
            \midrule
            Mamba-2 & 16.56 & 29.8 & 13.5 \\
            DeltaNet & 17.72 & 26.2 & 13.6 \\
            Gated DeltaNet & 16.42 & 30.6 & 16.6 \\
            \rowcolor{KAUSTcyan!12} + sliding-window attn. & 16.07 & 40.1 & 17.8 \\
            \bottomrule
        \end{tabular}\end{adjustbox}

        \vspace{0.5em}
        \scriptsize Recall: average of recall-intensive tasks. Long: LongBench average.
    \end{columns}
    \vspace{0.3em}
    \src{Yang, Kautz, Hatamizadeh, ``Gated Delta Networks,'' ICLR 2025, \href{https://arxiv.org/abs/2412.06464}{arXiv:2412.06464}. Table as compiled in Yang's DeltaNet talk (2025), slide 27. Kimi Team, \href{https://arxiv.org/abs/2510.26692}{arXiv:2510.26692}.}
\end{frame}

\begin{frame}{Making It Fast: Chunkwise Training}
    \begin{itemize}\itemsep0.3em
        \item A recurrence that is cheap on paper can still be slow on a GPU. Step-by-step updates use no matrix multiplies, so the tensor cores sit idle.
    \end{itemize}
    \vspace{0.2em}
    \begin{center}
    \small
    \begin{adjustbox}{max width=\linewidth}\begin{tabular}{llll}
        \toprule
        \textbf{Form} & \textbf{Sequential steps} & \textbf{Cost} & \textbf{Used for} \\
        \midrule
        Recurrent & $L$ & $O(Ld^2)$ & decoding \\
        Fully parallel & 1 & $O(L^2 d)$ & short sequences \\
        \rowcolor{KAUSTcyan!12} Chunkwise, chunk $C$ & $L/C$ & $O(LCd + Ld^2)$ & training \\
        \bottomrule
    \end{tabular}\end{adjustbox}
    \end{center}
    \vspace{0.2em}
    \begin{itemize}\itemsep0.4em
        \item Inside a chunk: ordinary masked attention. Across chunks: pass the state. In practice $C$ is 64 to 256.
        \item More FLOPs than the recurrence, yet faster in wall-clock time. \textbf{Hardware use beats FLOP count.}
        \item The delta rule needs one more trick: its chain of erase matrices is kept in a compact form (the WY representation).
    \end{itemize}
    \src{Yang et al., \href{https://arxiv.org/abs/2312.06635}{arXiv:2312.06635} and \href{https://arxiv.org/abs/2406.06484}{arXiv:2406.06484}. Library: \url{https://github.com/fla-org/flash-linear-attention}.}
\end{frame}
